% Volume VII: Formal Synthesis
% Part XIII. A Mathematics of Disciplined Distrust
% Chapter 75: Privacy–Publicity Duality
% Spine: proof-spines/ch075-spine.md

\chapter{Privacy--Publicity Duality}
\label{ch:ch075}

This chapter treats \textbf{privacy--publicity duality} as a joint constraint on legitimate institutional life. Let \(x\) denote personal information and \(g\) institutional action. Privacy constrains observation:
\[
\mathcal O(x)\leq \kappa.
\]
\ToyModel\ Privacy constrains observation by limiting personal observability.
Publicity constrains hidden authorization:
\[
\operatorname{Opacity}(g)\leq \rho.
\]
\ToyModel\ The publicity constraint limits how much official action may remain shielded from contestation.
The first inequality limits how much of a person's life may be rendered observable; the second limits how much official action may remain shielded from contestation.

A legitimate order must satisfy both constraints together. Maximum privacy can degenerate into concealed private domination, while maximum publicity can degenerate into universal exposure. The chapter's core claim is therefore dual rather than scalar: freedom occupies neither endpoint. Reliable evidence depends on protected zones in which appearances do not instantly harden into public dossier, and legitimate power depends on publicity sufficient to prevent hidden authorization from becoming unanswerable rule.
